% ============================================================
% chapters/02_etat_art.tex //  : solutions existantes et positionnement d'AGT Infra
% ============================================================
\chapter{État de l'art}
\label{chap:etat-art}



\agtlettrine{A}vant de détailler la conception d'AGT Infra, il est légitime de se demander
pourquoi ce projet n'a pas simplement reposé sur un outil déjà existant sur
le marché. Ce chapitre répond à cette question de manière factuelle : il
présente les principales solutions utilisées aujourd'hui pour gérer des
serveurs et des conteneurs Docker, puis explique en quoi leur périmètre ne
correspondait pas exactement au besoin d'AG Technologies. L'objectif n'est
pas de dénigrer ces outils, largement utilisés et éprouvés en production,
mais de justifier un choix d'ingénierie à partir de critères objectifs.
\minisommaire
\section{Critères de comparaison retenus}

Ces critères ont été définis à partir du contexte réel d'AG Technologies
(chapitre \ref{chap:entreprise}), avant toute comparaison, afin d'éviter de
choisir des critères qui favoriseraient artificiellement AGT Infra.

\begin{itemize}
  \item \textbf{Gestion multi-serveurs sans agent lourd} : capacité à piloter
  plusieurs serveurs physiques ou VPS uniquement via SSH, sans installer de
  démon ou d'agent supplémentaire exposé sur le réseau.
  \item \textbf{Orchestration de conteneurs Docker existants} : capacité à
  gérer des conteneurs Docker déjà en place, sans imposer une nouvelle couche
  de virtualisation ou d'exécution.
  \item \textbf{Contrôle d'accès granulaire et évolutif} : possibilité de
  définir des rôles et des permissions fins (par ressource, par action), et
  de les modifier sans redéployer l'outil.
  \item \textbf{Supervision temps réel intégrée} : suivi des métriques
  système (CPU, RAM, disque) avec des seuils d'alerte personnalisables.
  \item \textbf{Console web intégrée} : accès direct à un terminal serveur
  depuis l'interface, sans ouvrir une session SSH séparée.
  \item \textbf{Coût} : licence, abonnement, ou infrastructure additionnelle
  nécessaire pour obtenir les fonctionnalités recherchées.
  \item \textbf{Maturité et activité du projet} : fréquence des mises à jour,
  état réel de la maintenance.
  \item \textbf{Cohabitation sur une infrastructure mutualisée} : capacité à
  cohabiter proprement avec d'autres applications sur un même serveur, sans
  nécessiter un cluster dédié.
\end{itemize}

\section{Panorama des solutions existantes}

\subsection{Portainer}

Portainer est une interface graphique de gestion Docker très répandue,
disponible en version communautaire gratuite (Community Edition) couvrant
Docker, Swarm, Kubernetes et Azure Container Instances depuis un même
tableau de bord \cite{portainer_setup_2026}. Cette capacité à piloter
plusieurs environnements Docker depuis un point unique est justement la
raison principale pour laquelle de nombreuses équipes l'adoptent
\cite{portainer_setup_2026}. En revanche, le contrôle d'accès fin (rôles
par environnement comme administrateur, opérateur, ou lecteur seul) fait
partie de l'édition payante Business Edition, pas de la version gratuite
\cite{portainer_rbac_2026, portainer_security_2026}. Portainer reste par
ailleurs une cible de choix pour les attaquants puisqu'il pilote directement
Docker : une alerte de sécurité belge a signalé en 2026 des failles
critiques touchant plusieurs versions du produit, rappelant l'importance
de le maintenir à jour et de ne jamais l'exposer sans protection
supplémentaire \cite{portainer_setup_2026}.

\subsection{Rancher}

Rancher est conçu pour administrer plusieurs clusters Kubernetes depuis une
seule interface, avec un RBAC fin appliqué à la fois au niveau du cluster
et au niveau du projet (équivalent d'un espace de noms Kubernetes)
\cite{rancher_zesty}. Son modèle de permissions est puissant mais complexe :
SUSE, l'éditeur du produit, a dû retravailler en profondeur le moteur RBAC
de Rancher en 2025 pour améliorer ses performances face à la multiplication
des liaisons de rôles dans les grands parcs de clusters
\cite{rancher_suse_rbac}. Cette complexité est justifiée quand l'entreprise
gère réellement des dizaines de clusters Kubernetes, mais représente un
investissement disproportionné pour un parc de VPS gérés directement via
Docker, sans orchestrateur Kubernetes sous-jacent.

\subsection{Proxmox VE}

Proxmox VE est avant tout un hyperviseur : il gère des machines virtuelles
et des conteneurs système LXC sur du matériel physique
\cite{proxmox_petronella}. Ce n'est qu'en 2026, avec les versions 8.3 et 9,
que Proxmox a commencé à expérimenter l'exécution d'images au format OCI
(le format utilisé par Docker) directement dans ses conteneurs LXC, sans
passer par un démon Docker complet \cite{proxmox_techfuel}. Cette fonction
reste toutefois qualifiée de version d'évaluation technologique par
l'éditeur lui-même : il n'existe encore aucun équivalent de Docker Compose
ni d'écosystème d'orchestration autour de cette approche
\cite{proxmox_techfuel, proxmox_mrplanb}. Proxmox répond donc à un besoin
de virtualisation d'infrastructure, différent du besoin d'AGT Infra qui est
de piloter des conteneurs Docker déjà déployés sur des serveurs existants.

\subsection{Cockpit}

Cockpit est une console d'administration web légère, intégrée par défaut à
plusieurs distributions Linux. Contrairement à une idée reçue, Cockpit sait
gérer plusieurs serveurs depuis un même tableau de bord : il suffit
d'ajouter un serveur distant, et la connexion s'établit ensuite via SSH
\cite{cockpit_redhat, cockpit_rhel}. Sa limite pour le contexte d'AGT Infra
se situe ailleurs : le contrôle d'accès repose entièrement sur les comptes
système du serveur (utilisateurs Linux, droits sudo), et non sur un modèle
de rôles et de permissions propre à l'application. Il n'existe donc pas de
notion de permission granulaire par ressource (par exemple, autoriser un
utilisateur à consulter la supervision d'un serveur sans lui donner accès
à sa console web), ni de pipeline de déploiement intégré.

\subsection{Ansible AWX et Semaphore}

AWX et Semaphore sont des interfaces web pour exécuter des playbooks
Ansible (des scripts d'automatisation de configuration), pas des outils de
supervision temps réel. AWX propose un RBAC très détaillé, avec des
permissions par objet (inventaire, modèle de tâche, identifiants) et une
synchronisation possible avec un annuaire d'entreprise
\cite{semaphore_awx_2026}. Ce niveau de détail a toutefois un coût : le
développement d'AWX est à l'arrêt depuis juillet 2024, l'équipe ayant
reconnu que l'architecture actuelle du projet limite sa capacité à évoluer,
et qu'une refonte est en cours sans calendrier annoncé
\cite{semaphore_awx_2026, rogerperkin_awx}. Semaphore, plus récent et plus
activement maintenu, propose un modèle de rôles plus simple, limité à
quatre rôles par projet (Owner, Manager, Task Runner, Guest), les
permissions personnalisées plus fines étant réservées à une offre
Entreprise payante \cite{semaphore_vs_awx, rogerperkin_awx}. Dans les deux
cas, ces outils automatisent des tâches ponctuelles, ils ne fournissent pas
de tableau de bord de supervision continue ni de console web intégrée.

\subsection{Coolify}

Coolify est une plateforme auto-hébergée de type PaaS (Platform as a
Service) : elle automatise le déploiement d'applications à partir d'un
dépôt Git, avec certificats SSL et sauvegardes programmées inclus
\cite{coolify_github}. Son modèle de permissions est volontairement simple,
limité à deux rôles (Administrateur ou Membre), et est explicitement
documenté comme n'étant pas adapté à une logique stricte de moindre
privilège : il n'existe pas de permission par service ou par action
\cite{coolify_admin_it}. Coolify effectue par nature des opérations
privilégiées sur le serveur hôte (déploiement, gestion Docker), ce qui en
fait une cible sensible : plusieurs failles critiques (score de gravité
maximal) ont été corrigées début 2026, liées à un défaut d'assainissement
des commandes système générées automatiquement
\cite{coolify_cve_2026}. Ce point rejoint une décision déjà actée dans
AGT Infra (règle numéro 6) : le backend doit rester l'unique garde-fou de
sécurité, précisément pour éviter ce type de surface d'attaque.

\section{Tableau comparatif synthétique}

% Colonnes a largeur fixe (p{}) plutot que tabularx : garantit qu'aucun
% mot ne peut deborder ou chevaucher la colonne voisine, meme sans point
% de cesure disponible (probleme rencontre avec des tokens comme
% "Gratuit/payant" en colonnes trop etroites).
\begin{table}[H]
\centering
\small
\renewcommand{\arraystretch}{1.3}
\begin{tabular}{p{3.4cm}*{7}{>{\centering\arraybackslash}p{1.3cm}}}
\toprule
\textbf{Critère} & \textbf{AGT} & \textbf{Portai.} & \textbf{Ranch.} & \textbf{Proxm.} & \textbf{Cockpit} & \textbf{AWX/Sem.} & \textbf{Coolify} \\
\midrule
Multi-serveurs (SSH seul)        & O & O   & N   & N & O & P   & N \\
Orchestration Docker existant    & O & O   & N   & N & P & N   & O \\
RBAC/IBAC fin et dynamique       & O & € & O   & N & N & P/€ & N \\
Supervision temps réel           & O & N   & P   & P & P & N   & N \\
Console web intégrée             & O & O   & O   & O & O & N   & N \\
Coût                             & Interne & 0/€ & 0/€ & 0 & 0 & 0/€ & 0/€ \\
Développement actif              & O & O   & O   & O & O & V   & O \\
Cohabitation VPS mutualisé       & O & P   & N   & N & P & P   & P \\
\bottomrule
\end{tabular}
\caption{Comparaison synthétique des solutions existantes face aux besoins d'AG Technologies}
\label{tab:comparatif-etat-art}
\end{table}

\begin{flushleft}
\footnotesize
\textbf{Légende} : O = Oui, N = Non, P = Partiel, V = Variable selon le
sous-projet (Semaphore actif, AWX en pause), 0 = Gratuit ou open-source,
€ = fonctionnalité payante requise pour couvrir ce critère.
\end{flushleft}

\section{Synthèse et justification du développement interne}

Cette comparaison ne conduit pas à la conclusion que les solutions
existantes seraient de mauvaise qualité, bien au contraire : Portainer,
Rancher, Proxmox, Cockpit, AWX/Semaphore et Coolify sont des projets
matures, largement adoptés, et bénéficient d'années de tests en production
que ne possède évidemment pas un outil interne récemment développé. Construire un outil en
interne signifie aussi en assumer seul la maintenance et la sécurité dans
la durée, une responsabilité qu'un outil du marché, déjà éprouvé, aurait
partiellement absorbée.

Ce que révèle néanmoins le tableau comparatif, c'est qu'aucune de ces
solutions ne réunit simultanément l'ensemble des critères retenus au
paragraphe 3.2. Chaque outil couvre bien une partie du besoin, mais avec un
angle différent : Portainer et Coolify ciblent la gestion ou le déploiement
de conteneurs sans contrôle d'accès fin gratuit ; Rancher offre un RBAC
puissant mais conçu pour Kubernetes, disproportionné pour un parc de VPS
sous Docker simple ; Proxmox reste avant tout un hyperviseur de machines
virtuelles ; Cockpit couvre bien plusieurs serveurs mais sans notion
applicative de rôles ; AWX et Semaphore automatisent des tâches
ponctuelles sans assurer de supervision continue.

Le besoin d'AG Technologies combine plusieurs contraintes en même temps :
un VPS partagé entre plusieurs applications internes (donc pas de cluster
dédié), un contrôle d'accès qui doit pouvoir évoluer sans redéploiement
(module RBAC/IBAC dynamique), une supervision continue avec des seuils
personnalisables, et une console web intégrée, le tout en gardant une
maîtrise complète du code dans un contexte où l'outil gère des accès
privilégiés à des serveurs de production. C'est cette combinaison précise
de contraintes, plus que l'absence d'alternative pour chaque besoin pris
isolément, qui justifie le développement d'AGT Infra en interne.